\section{Your independent capstone}
Choose one thread and keep the first deliverable narrow enough that every assertion can be checked. An assignments project can add integer capacities to tasks and derive the appropriate shortage condition by replacing a capacity-$c$ task with $c$ distinct slots. An iteration project can analyze a specified decaying-error schedule. A sign-sequence project can give a reproducible classification at one additional finite length and investigate which symmetries preserve its conditions. These are starting points, not claims of new results.

Submit a precise statement, a prerequisite map, a proof or counterexample, executable checks when relevant, an assumption-and-evidence record, and a short explanation of what is genuinely new to the literature versus new only to you. The oral component should ask you to explain one delicate step, change an assumption, and predict what fails. Your ability to answer matters more than the fluency of the saved AI conversation.

An effective final review uses the assistant in three passes with different tasks. First request a faithful reconstruction of the argument. Next request counterexamples to stronger variants and scrutiny of each hypothesis. Finally ask for a correspondence check between the prose, formulas, code, and formal statements. These passes may still share blind spots; they do not constitute independent human peer review. Their value lies in the concrete obligations they expose for you to check.

\begin{lab}{From a plausible assertion to a defensible result}
Start with the false claim ``bounded update errors do not affect the limiting fixed point.'' Produce the constant-bias counterexample before asking for a repair. Derive the one-step inequality yourself, ask the assistant to propose the unfolded form, and verify its indexing at $n=1,2,3$. Then prove the induction step, test the sharp example, and explain why alternating bounded errors can prevent convergence. Finally close the chapter and reconstruct the finite-prefix argument for vanishing errors. This completes the collaboration loop: you can now explain the result rather than merely possess its generated proof.
\end{lab}
